\documentclass[12pt]{article}
\usepackage[T1]{fontenc}
\usepackage[utf8]{inputenc}
\usepackage{lmodern}
\usepackage{textcomp}
\usepackage{enumitem}
\usepackage{geometry}
\usepackage{graphicx}
\usepackage{amsmath, amssymb}
\geometry{margin=1in}
\begin{document}
\begin{center}
Chemistry - Undergraduate Level Assessment 3 \
Total Marks: 40 \
Answer all questions.
\end{center}

\section*{Multiple Choice (10 marks)}
\begin{enumerate}[label=\arabic*.]
\item Which of the following lists the hydrides of group-14 elements in order of thermal stability, from lowest to highest?
\begin{enumerate}[label=(\alph*)]
    \item PbH 4 \textless{} SnH 4 \textless{} GeH 4 \textless{} SiH 4 \textless{} CH 4
    \item PbH 4 \textless{} SnH 4 \textless{} CH 4 \textless{} GeH 4 \textless{} SiH 4
    \item CH 4 \textless{} SiH 4 \textless{} GeH 4 \textless{} SnH 4 \textless{} PbH 4
    \item CH 4 \textless{} PbH 4 \textless{} GeH 4 \textless{} SnH 4 \textless{} SiH 4
\end{enumerate}
\item What is the limiting high-temperature molar heat capacity at constant volume (C\_V) of a gas-phase diatomic molecule?
\begin{enumerate}[label=(\alph*)]
    \item 1.5 R
    \item 2 R
    \item 2.5 R
    \item 3.5 R
\end{enumerate}
\item The T$_{2}$ of an NMR line is 15 ms. Ignoring other sources of linebroadening, calculate its linewidth.
\begin{enumerate}[label=(\alph*)]
    \item 0.0471 Hz
    \item 21.2 Hz
    \item 42.4 Hz
    \item 66.7 Hz
\end{enumerate}
\item The normal modes of a carbon dioxide molecule that are infrared-active include which of the following? I. Bending II. Symmetric stretching III. Asymmetric stretching
\begin{enumerate}[label=(\alph*)]
    \item I only
    \item II only
    \item III only
    \item I and III only
\end{enumerate}
\item Considering 0.1 M aqueous solutions of each of the following, which solution has the lowest pH?
\begin{enumerate}[label=(\alph*)]
    \item Na$_{2}$CO$_{3}$
    \item Na$_{3}$PO$_{4}$
    \item Na$_{2}$S
    \item NaCl
\end{enumerate}
\end{enumerate}

\section*{True / False (10 marks)}
\textit{Answer True or False}
\begin{enumerate}[label=\arabic*.]
\setcounter{enumi}{5}
\item True or False: Protons in a molecule with a rotational correlation time of 1 ns have the fastest spin-lattice relaxation at a magnetic field of 3.74 T.
\item True or False: The nuclide 19 F has an NMR frequency of 115.5 MHz in a 20.0 T magnetic field.
\item True or False: AgCl has a lower melting point than HCl because it is an ionic compound with strong ionic bonds.
\item True or False: Most organic compounds do not contain silicon atoms, making tetramethylsilane an ideal reference for 1 H NMR spectroscopy.
\item True or False: The molecular geometry of thionyl chloride, SOCl 2, is tetrahedral because it has four bonding pairs of electrons.
\end{enumerate}

\section*{Long Form Response (20 marks)}
\textit{Answer each question with a clear and complete explanation.}
\begin{enumerate}[label=\arabic*.]
\setcounter{enumi}{10}
\item Discuss the factors influencing the strength of a 13 C$_{90}$^\ensuremath{\circ} pulse of duration 1 μs, and calculate its value, providing a detailed explanation of your reasoning.
\item Discuss the concept of dissociation energy in the context of the hydrogen-bromine bond, specifically analyzing the enthalpy change associated with the reaction HBr(g) → H(g) + Br(g).
\end{enumerate}

\vfill
\noindent\textit{End of Paper}
\end{document}